Three hundred problems on the frontier of number theory, algebra and probability. How many choices of signs make
$\pm1\pm2\pm\dots\pm n$ vanish? Why does a random polynomial of degree one hundred have only about three real roots? Why does a
sum of random signs weighted by powers of the golden ratio have no density, although it looks perfectly smooth? Each problem
mixes the tools of the third part of \emph{Order and Chance}: sharp limit theorems, Brownian motion, algebraic number theory,
curves over finite fields, sieves, and the zeros of the zeta function seen through random matrices. None says which tool to
use: finding it is the game. Hints point the way when the search stalls, every solution begins with how one might discover
the approach before giving a complete proof, and long \emph{Journeys} walk step by step to famous theorems. Remarks show
where each problem meets current research and open questions. Companion Python notebooks check every number and turn the
ideas into experiments.
